% Tài liệu tham khảo và khai báo sử dụng công cụ AI
% Các mục Breiman, Brier, Efron, Fawcett, Flores, Hastie, Hubert, Pedregosa, Rousseeuw, Saito,
% Simon, Sprague, Turban-Sharda-Delen, Youden, Lemmens lấy từ danh sách "Tài liệu tham khảo chương 3" của nhóm trên Google Doc (tab t.0).
% Elmasri, Han, Inmon, Kimball, tài liệu môn học, Turban-Aronson có sẵn trong báo cáo LaTeX từ trước.
% Nielsen (1994), ISO 9241-11 và Streamlit được thêm khi viết Chương 6 và 7.
% Buckinx (2005) và Neslin (2006) thêm ngày 07/11 sau khi đối chiếu tóm tắt bài báo, bỏ nguồn web Empire Data Systems.
% Han 3rd ed. phát hành 2011 (bản in ghi bản quyền 2012). Turban-Sharda-Delen 9th ed. bản quyền 2011 (in 2010).
\addcontentsline{toc}{section}{Tài liệu tham khảo}
\begin{thebibliography}{99}
\bibitem{breiman2001} L. Breiman, ``Random forests'', \textit{Machine Learning}, 45(1), 2001, pp. 5--32.
\bibitem{brier1950} G. W. Brier, ``Verification of forecasts expressed in terms of probability'', \textit{Monthly Weather Review}, 78(1), 1950, pp. 1--3.
\bibitem{buckinx2005} W. Buckinx, D. Van den Poel, ``Customer base analysis: partial defection of behaviourally loyal clients in a non-contractual FMCG retail setting'', \textit{European Journal of Operational Research}, 164(1), 2005, pp. 252--268.
\bibitem{efron1993} B. Efron, R. J. Tibshirani, \textit{An Introduction to the Bootstrap}, Chapman \& Hall/CRC, 1993.
\bibitem{elmasri2016} R. Elmasri, S. B. Navathe, \textit{Fundamentals of Database Systems}, 7th ed., Addison-Wesley, 2016.
\bibitem{fawcett2006} T. Fawcett, ``An introduction to ROC analysis'', \textit{Pattern Recognition Letters}, 27(8), 2006, pp. 861--874.
\bibitem{flores1986} B. E. Flores, D. C. Whybark, ``Multiple criteria ABC analysis'', \textit{International Journal of Operations \& Production Management}, 6(3), 1986, pp. 38--46.
\bibitem{han2011} J. Han, M. Kamber, J. Pei, \textit{Data Mining: Concepts and Techniques}, 3rd ed., Morgan Kaufmann, 2011.
\bibitem{hastie2009} T. Hastie, R. Tibshirani, J. Friedman, \textit{The Elements of Statistical Learning: Data Mining, Inference, and Prediction}, 2nd ed., Springer, 2009.
\bibitem{hubert1985} L. Hubert, P. Arabie, ``Comparing partitions'', \textit{Journal of Classification}, 2(1), 1985, pp. 193--218.
\bibitem{inmon2005} W. H. Inmon, \textit{Building the Data Warehouse}, 4th ed., Wiley, 2005.
\bibitem{iso9241} ISO 9241-11:2018, \textit{Ergonomics of human-system interaction --- Part 11: Usability: Definitions and concepts}, International Organization for Standardization, 2018.
\bibitem{kimball2013} R. Kimball, M. Ross, \textit{The Data Warehouse Toolkit: The Definitive Guide to Dimensional Modeling}, 3rd ed., Wiley, 2013.
\bibitem{lemmens2020} A. Lemmens, S. Gupta, ``Managing churn to maximize profits'', \textit{Marketing Science}, 39(5), 2020, pp. 956--973.
\bibitem{neslin2006} S. A. Neslin, S. Gupta, W. Kamakura, J. Lu, C. H. Mason, ``Defection detection: measuring and understanding the predictive accuracy of customer churn models'', \textit{Journal of Marketing Research}, 43(2), 2006, pp. 204--211.
\bibitem{nielsen1994} J. Nielsen, ``Enhancing the explanatory power of usability heuristics'', \textit{Proceedings of the SIGCHI Conference on Human Factors in Computing Systems (CHI '94)}, ACM, 1994, pp. 152--158.
\bibitem{pedregosa2011} F. Pedregosa et al., ``Scikit-learn: machine learning in Python'', \textit{Journal of Machine Learning Research}, 12, 2011, pp. 2825--2830, \url{https://scikit-learn.org/stable/user_guide.html}.
\bibitem{tlmh} Phan Trọng Nhân, \textit{Data Warehouses and Decision Support Systems}, tài liệu môn học, Trường Đại học Bách Khoa, ĐHQG-HCM, 2026.
\bibitem{rousseeuw1987} P. J. Rousseeuw, ``Silhouettes: a graphical aid to the interpretation and validation of cluster analysis'', \textit{Journal of Computational and Applied Mathematics}, 20, 1987, pp. 53--65.
\bibitem{saito2015} T. Saito, M. Rehmsmeier, ``The precision-recall plot is more informative than the ROC plot when evaluating binary classifiers on imbalanced datasets'', \textit{PLoS ONE}, 10(3), 2015, e0118432.
\bibitem{simon1977} H. A. Simon, \textit{The New Science of Management Decision}, revised ed., Prentice Hall, 1977.
\bibitem{sprague1980} R. H. Sprague, ``A framework for the development of decision support systems'', \textit{MIS Quarterly}, 4(4), 1980, pp. 1--26.
\bibitem{streamlit} Streamlit Inc., \textit{Streamlit documentation}, \url{https://docs.streamlit.io}, truy cập tháng 11/2026.
\bibitem{turban2005} E. Turban, J. E. Aronson, \textit{Decision Support Systems and Intelligent Systems}, 7th ed., Prentice Hall, 2005.
\bibitem{turban2011} E. Turban, R. Sharda, D. Delen, \textit{Decision Support and Business Intelligence Systems}, 9th ed., Pearson, 2011.
\bibitem{youden1950} W. J. Youden, ``Index for rating diagnostic tests'', \textit{Cancer}, 3(1), 1950, pp. 32--35.
\end{thebibliography}

\vspace{1.2em}
\section*{Khai báo sử dụng công cụ AI}
\addcontentsline{toc}{section}{Khai báo sử dụng công cụ AI}

Nhóm có sử dụng công cụ AI để hỗ trợ làm báo cáo, theo quy định của môn học là chỉ dùng để hỗ trợ và phải khai báo công cụ cùng cách dùng. Bảng~\ref{tab:khaibaoai} nêu từng công cụ và việc nó được dùng vào.

\begin{table}[!htbp]
\centering
\caption{Công cụ AI đã sử dụng và mục đích sử dụng}
\label{tab:khaibaoai}
\footnotesize
\begin{tabular}{p{3.5cm}p{11.5cm}}
\toprule
Công cụ & Cách sử dụng \\
\midrule
NotebookLM (Google) & Nạp các slide môn học làm nguồn, rồi dùng để tra cứu nội dung lý thuyết trong slide, đối chiếu từng câu lý thuyết của báo cáo với slide trước khi đưa vào bài (chỉ ghi ``theo tài liệu môn học'' khi slide có nội dung đó), và phản biện cách trình bày. \\
\mbox{Claude (Anthropic)} & Hỗ trợ soạn thảo và chỉnh sửa nội dung báo cáo, chuyển nội dung sang LaTeX, viết các đoạn mã Python sinh bảng và hình từ file kết quả, chạy lại và đối chiếu số liệu trong báo cáo với kết quả phân tích, chụp lại màn hình ứng dụng, biên dịch báo cáo và kiểm tra bố cục. \\
\bottomrule
\end{tabular}
\end{table}

Các công cụ này không thay nhóm trong phần chuyên môn. Dữ liệu được nhóm xử lý và các mô hình được nhóm huấn luyện bằng mã chạy trên dữ liệu của đề tài. Mọi con số trong báo cáo lấy từ kết quả chạy đó, và nội dung do AI gợi ý đều được nhóm đối chiếu lại với slide môn học hoặc với số liệu trước khi đưa vào báo cáo. Nhóm chịu trách nhiệm về toàn bộ nội dung báo cáo.
